\documentclass[11pt]{article}
\usepackage[T1]{fontenc}
\usepackage[cmex10]{amsmath}
\usepackage{amssymb}
\usepackage[margin=1in]{geometry}
\usepackage[hidelinks]{hyperref}
\DeclareFontShape{OMX}{cmex}{m}{n}{<->cmex10}{}
\hypersetup{pdftitle={The sharp lower zero edge for ordered partitions counted by large blocks},pdfauthor={ProgramComputer}}
\title{The sharp lower zero edge for ordered partitions counted by large blocks}
\author{ProgramComputer}
\date{November 25, 2026}
\begin{document}
\maketitle
\begin{abstract}
For each integer threshold $j\ge2$, we determine the infimum of the absolute values of the negative real zeros of the ordered-partition polynomials counting blocks of size at least $j$. Every row is strictly positive up to this boundary, which is described by a unique double pole. A positive-real-part estimate and a Fej\'er kernel show that the critical pole is isolated in modulus for every threshold. On an interval just past the boundary, a fraction bounded away from zero of every sufficiently large row consists of simple negative zeros strictly alternating with the next row. The local limit is a sinc function. Its derivatives give Bessel constants for the innermost real zeros of all fixed derivatives of the enumerating polynomial.
\end{abstract}
\section{A sharp lower bound for real root moduli at every ordered threshold}

For an integer $j\ge2$, let $O_{n,j}(t)$ count ordered set partitions of $[n]$ by their blocks of size at least $j$. Define
\[
S_{j-1}(z)=\sum_{s=0}^{j-1}\frac{z^s}{s!},\qquad
T_j(z)=e^z-S_{j-1}(z),\qquad
D_{j,q}(z)=2-S_{j-1}(z)+qT_j(z).
\]
The ordered sequence construction gives $\sum_{n\ge0}O_{n,j}(-q)z^n/n!=1/D_{j,q}(z)$. The question concerns the lower edge of the moduli of negative real zeros. At thresholds two and three, all roots in every row are known to interlace by the companion arguments. See~\cite{original,companion}.

\textbf{Theorem.} Let $\rho_j$ be the unique positive solution of
\[
e^{\rho_j}-2=2\,\frac{T_j(\rho_j)}{\rho_j^{j-1}/(j-1)!},
\qquad
 q_j^*=\frac{S_{j-1}(\rho_j)-2}{T_j(\rho_j)}.
\]
Then $q_j^*>0$ and $O_{n,j}(-q)>0$ for every $n\ge0$ and $0\le q\le q_j^*$. At $q=q_j^*$, $D_{j,q}$ has a double zero at $\rho_j$, and every other zero has modulus strictly greater than $\rho_j$. For each $q>q_j^*$ there are infinitely many $n$ with $O_{n,j}(-q)<0$. In particular, $q_j^*$ is the infimum of the moduli of the negative real zeros over all rows. Every neighborhood of this edge above $q_j^*$ contains a zero from infinitely many rows.

\section{The critical radius}

For $r>0$, write
\[
A_j(r)=\frac{T_j(r)}{r^{j-1}/(j-1)!}
      =\sum_{\ell\ge1}\frac{r^\ell}{j(j+1)\cdots(j+\ell-1)},
\qquad E_j(r)=e^r-2-2A_j(r).
\]
The coefficient of $r^\ell$ in $E_j$ is nonnegative for $\ell\ge1$, since $j(j+1)\cdots(j+\ell-1)\ge(\ell+1)!$. It is positive for $\ell\ge2$. Thus $E_j$ increases strictly from $-1$ to infinity and has exactly one positive zero. At fixed $r>0$, $A_j(r)$ decreases strictly as $j$ increases. Consequently $\rho_j$ decreases with $j$.

A positive zero of $D_{j,q}$ at $r$ requires $q=Q_j(r)=(S_{j-1}(r)-2)/T_j(r)$. The derivative of $Q_j$ has the sign of $-E_j(r)$. In fact, writing $w=r^{j-1}/(j-1)!$, the numerator of its derivative is
\[
S_{j-2}(r)T_j(r)-(S_{j-1}(r)-2)(T_j(r)+w)
=-wE_j(r).
\]
The ratio $Q_j$ is zero where $S_{j-1}=2$, has positive derivative there, and tends to zero at infinity. It has a strict maximum $q_j^*>0$ at $\rho_j$ with negative second derivative. The identity $D_{j,q}(r)=T_j(r)(q-Q_j(r))$ shows that this critical zero is double. It also shows that $D_{j,q}$ has no positive zero when $q>q_j^*$.

The critical equation also gives $q_j^*=2(j-1)!/\rho_j^{j-1}-1$. For $\lambda=\log 2$ and large $j$, the series for $A_j$ yields $A_j(\rho_j)=\rho_j/j+O(j^{-2})$, uniformly because the radii are bounded. Using $e^{\rho_j-\lambda}=1+A_j(\rho_j)$ gives
\[
\rho_j=\lambda\left(1+\frac1j\right)+O(j^{-2}),
\qquad q_j^*\sim \frac{2(j-1)!}{e\,\lambda^{j-1}}.
\]

\section{Positivity at the critical value}

Put $\rho=\rho_j$, $q=q_j^*$, and $D(z)=D_{j,q}(\rho z)$. First suppose $\rho\ge1$. The coefficient of $z^k$ in $G(z)=D(z)/(1-z)$ is the sum of the first $k+1$ coefficients of $D$. For $1\le k<j$, this is $1-\sum_{s=1}^k\rho^s/s!\le0$. For $k\ge j$, it is $-q\sum_{s>k}\rho^s/s!<0$, by $D(1)=0$. Hence $B=1-G$ has nonnegative coefficients and zero constant term. The expansion $(1-z)^{-1}\sum_{k\ge0}B(z)^k$ proves strict coefficient positivity of $1/D$. At the double zero, $G(1)=0$, so $B(1)=1$. For $|z|<1$, $|B(z)|<1$. On $|z|=1$, equality $B(z)=1$ forces all powers with positive coefficients to equal $1$. Their support contains consecutive integers, so this occurs only at $z=1$. Thus the critical zero is isolated strictly in modulus in this case.

For the remaining radii use $H(z)=D(z)/(1-z)^2$, an entire function. Write $a_s=\rho^s/s!$ and $b_s=q\rho^s/s!$. For $z=e^{i\theta}$ with $x=1-\cos\theta>0$, the identity $(1-z)^2=-2xz$ and $D(1)=0$ give
\[
2\Re H(e^{i\theta})=1-\sum_{s=2}^{j-1}a_s R_{s-1}(\theta)+\sum_{s\ge j}b_s R_{s-1}(\theta),
\qquad R_k(\theta)=\frac{1-\cos(k\theta)}{1-\cos\theta}.
\]
The quotient has a continuous extension at zero with value $k^2$. Also $0\le R_k\le k^2$, by the geometric-sum expression for $|1-e^{ik\theta}|/|1-e^{i\theta}|$.

The elementary bounds on the critical radii used below are
\[
\rho_3>1,\qquad \frac{83}{100}<\rho_6,\qquad
\rho_4<\frac{19}{20},\quad \rho_5<\frac{89}{100},\quad
\rho_6<\frac{17}{20},\quad \rho_7<\frac56.
\]
They follow by evaluating the sign of $E_j$ at the given rational arguments. For $0<x\le1$, bound $e^x$ below by $\sum_{s=0}^6 x^s/s!$ and above by this sum plus $x^7/(7!(1-x/8))$. Bound $A_j(x)$ below by its first two terms and above by $x/j+x^2/(j(j+1)(1-x/(j+2)))$. The rational comparisons are verified in the companion script. In particular, the preceding $G$ argument applies at $j=2,3$.

For $j\ge7$, monotonicity gives $\rho<5/6$, and the sum of all the negative contributions is bounded above by
\[
\sum_{s\ge2}(s-1)^2\frac{(5/6)^s}{s!}
=\frac{31}{36}e^{5/6}-1<1.
\]
For $j=4,5,6$, let $U_j$ be the indicated rational upper bound on $\rho_j$. If $x\ge1/4$, then $R_2=4-2x\le7/2$ and $R_3=(3-2x)^2\le25/4$. Together with $R_1=1$ and $R_4\le16$ these give the first column of strict upper bounds below. For $0<x\le1/4$ we may take $\theta\in[0,\pi]$ by symmetry and have $\theta\le\pi/3$ and $\theta\le(j-1)\theta\le2\pi-\theta$. Therefore $R_{j-1}\ge1$. Since
\[
b_j=(S_{j-1}(\rho)-2)\frac{\rho^j/j!}{T_j(\rho)}
>\frac45\left(\sum_{s=0}^3\frac{(83/100)^s}{s!}-2\right)=:L,
\]
the $s=j$ contribution supplies at least $L$. Here the tail ratio bound gives $(\rho^j/j!)/T_j(\rho)\ge1-\rho/(j+1)\ge j/(j+1)\ge4/5$. All arithmetic in the following upper bounds is rational. The decimal values are rounded up.

\[
\begin{array}{c|cc}
 j & x\ge1/4 & x\le1/4\text{, after subtracting }L\\ \hline
 4 & 0.952 & 0.808\\
 5 & 0.971 & 0.886\\
 6 & 0.915 & 0.810
\end{array}
\]

Thus $\Re H$ is strictly positive on the unit circle, including at $1$ by continuity and the strict uniform bounds. The minimum principle for harmonic functions gives $\Re H>0$ on the closed disk. This also rules out every other zero of $D$ there. The reciprocal $K=1/H$ is analytic on a neighborhood of the closed disk and has positive real part there. Its real Taylor coefficients $K_k$ give
\[
[z^n]\frac{1}{D(z)}
=\sum_{k=0}^n(n+1-k)K_k
=\frac{1}{2\pi}\int_{-\pi}^{\pi}\Re K(e^{i\theta})
  \left|\sum_{k=0}^n e^{\mathrm{i}k\theta}\right|^2 d\theta>0.
\]
The integral identity follows by expanding the square, whose Fourier coefficients are $n+1-|k|$ for $|k|\le n$. The integrand is positive away from the finitely many zeros of the squared sum. Rescaling $z$ by the positive $\rho$ preserves coefficient signs.

\section{All subcritical parameters and sharpness}

Write $F_p(z)=1/D_{j,p}(z)$. For $0\le q\le q_j^*$ the exact identity
\[
F_q=\frac{F_{q_j^*}}{1-(q_j^*-q)T_jF_{q_j^*}}
\]
has a coefficientwise nonnegative geometric expansion, with the strictly positive coefficients of $F_{q_j^*}$ as its first term. It proves the claimed positivity throughout the interval.

For $q>q_j^*$, the reciprocal $F_q$ is analytic across the whole positive real axis and tends to zero there at infinity. Suppose its coefficients were eventually nonnegative. Add a polynomial to make the full series nonnegative. A nonnegative Taylor series of finite radius $R>0$ must be singular at $R$. Indeed, if it extended analytically across $R$, each Taylor coefficient of its expansion about $R$ would be the limit of a nonnegative derivative series from below. Expanding at $R+\varepsilon$ and interchanging the nonnegative summations would give convergence of the original series at $R+\varepsilon$, a contradiction. Here there is no positive real singularity, so the series after adding the polynomial would be entire. It would have infinitely many nonzero coefficients, since $F_q(0)=1$ and $F_q(x)\to0$ imply $F_q$ is not a polynomial. Nonnegative coefficients with infinite support force growth along the positive real axis exceeding every fixed polynomial bound. That contradicts the behavior of $F_q$ plus the added polynomial. Thus infinitely many coefficients of $F_q$ are negative. For every corresponding row, continuity yields a real zero with modulus strictly between $q_j^*$ and $q$.

The result concerns $j\ge2$. At $j=1$, every block carries the variable and every positive-size row has a zero at the origin, so that threshold has a different endpoint behavior.

\section{Consecutive interlacing on an interval at every fixed threshold}

Fix a block threshold $j\ge2$. Write $D(z,q)=1-S(z)+qC(z)$ with $S(z)=\sum_{a=1}^{j-1}z^a/a!$ and $C(z)=\sum_{a\ge j}z^a/a!$. Let $q_*$ and $\rho>0$ be the critical parameter and its double zero from the preceding theorem. That theorem gives $O_{n,j}(-q)>0$ throughout $0\le q\le q_*$ for every $n$, and $\rho$ is the only zero of $D(\cdot,q_*)$ on or inside $|z|=\rho$.

\textbf{Theorem.} For every fixed $j$ there is a nonempty interval immediately to the right of $q_*$ on which all sufficiently large consecutive rows have strictly alternating simple real zeros in the variable $q$. The number of such zeros in each row grows linearly with $n$. More precisely there is an analytic phase $\phi(q)>0$ on that interval with $\phi(q_*)=0$ by continuity, and for each sufficiently small fixed $q_0>q_*$,
\[
\#\{q\in(0,q_0]:O_{n,j}(-q)=0\}=\frac{n\phi(q_0)}{\pi}+O(1).
\]
At each fixed index $h\ge1$ the $h$th positive real zero satisfies
\[
n^2(q_{n,h}-q_*)\longrightarrow
  \frac{\rho^2D_{zz}(\rho,q_*)}{2C(\rho)}\pi^2 h^2 .
\]
This theorem controls the real zeros at the inner edge. The companion outer-edge theorem~\cite{companion} controls largest moduli among all complex zeros.

\textbf{The two poles near the double zero.} At the fold, $D_q=C(\rho)>0$ and $D_{zz}(\rho,q_*)>0$ by the double-zero minimum. Solve $D(z,q)=0$ as $q=Q(z)$. Then
\[
Q(z)=q_*-\frac{D_{zz}(\rho,q_*)}{2C(\rho)}(z-\rho)^2+O((z-\rho)^3).
\]
Taking an analytic square root of the nonvanishing factor in $(q_*-Q(z))/(z-\rho)^2$ and inverting locally gives, at $q=q_*+t^2$, a branch $z(t)=\rho+i\sqrt{2C(\rho)/D_{zz}(\rho,q_*)}\,t+O(t^2)$. The other zero is $z(-t)$, the conjugate for real $t$. Write
\[
\phi(t)=\frac{\log z(t)-\log z(-t)}{2i},\qquad
 r(t)=\exp\big((\log z(t)+\log z(-t))/2\big),
\]
choosing the real logarithm at $\rho$. The first function is odd and has positive derivative $\kappa=\sqrt{2C(\rho)/D_{zz}(\rho,q_*)}/\rho$. Use $\phi$ itself as the analytic parameter. Then $q(\phi)=q_*+\phi^2/\kappa^2+O(\phi^4)$ is even and strictly increasing for small positive $\phi$. The two poles are $r(\phi)e^{\pm i\phi}$.

\textbf{Uniform coefficient estimate.} Isolation at the fold gives a circle $|z|=R>\rho$ containing just the double zero at $q=q_*$. For small complex $\phi$ there are just the two continuations inside and $D$ stays uniformly away from zero on that circle. Put $N=n+1$ and $f_n(q)=O_{n,j}(-q)/n!$. Cauchy's formula enlarged to this circle gives the negative of the two residues of $1/(z^N D)$, with an analytic remainder $J_n$ satisfying $|J_n|\le C_0 R^{-n}$. This remainder is even in $\phi$ by its contour representation.

Write $D_z(r(\phi)e^{i\phi},q(\phi))=i\phi d(\phi)$, where $d(0)=\rho D_{zz}(\rho,q_*)>0$. Define the even and odd analytic functions $a$ and $B$ by
\[
a(\phi)=\exp\big((\log d(\phi)+\log d(-\phi))/2\big),\qquad
 B(\phi)=(\log d(\phi)-\log d(-\phi))/(2i).
\]
On the real axis $a>0$ and $B$ is real. Substitution of the residues gives the signed identity
\[
\frac{\phi a(\phi)r(\phi)^N}{2} f_n(q(\phi))
 =\sin(N\phi+B(\phi))+E_n(\phi),\qquad
 E_n=\phi a(\phi)r(\phi)^N J_n/2 .
\]
It extends analytically across zero. On a fixed smaller complex disk, $|E_n(\phi)|\le C_1|\phi|\sigma^n$ and $|E_n'(\phi)|\le C_2\sigma^n$ for some $\sigma<1$. Indeed $|r|<R$ on the larger disk after shrinking, and the derivative estimate follows by Cauchy's inequality.

\textbf{Counting and consecutive order.} On a fixed small interval $0\le\phi\le2\eta$, the phase $N\phi+B(\phi)$ starts at zero and has derivative $N+O(1)$. Wherever its distance from $\pi\mathbb Z$ is at least $\pi/6$, the sine has modulus at least $1/2$. In each full band within $\pi/6$ of $h\pi$, the right-hand side is strictly monotone since the principal derivative has modulus at least $\sqrt3(N-O(1))/2$ and fixed sign. Its endpoint signs are opposite. The half-band at the origin has just its forced zero from the normalizing factor. Thus the roots in $0<\phi<\eta$ have consecutive indices starting at one, one per complete band, with at most one unfinished band at the counting endpoint. The phase-count formula follows, and each root is simple in $q$ by $q'(\phi)>0$.

At the old $h$th zero, write $U=N\phi+B(\phi)$. Then $|\sin U|\le C_1\phi\sigma^n$ and $\cos U$ has sign $(-1)^h$ and magnitude at least $\sqrt3/2$. In the next row the phase is $U+\phi$. Because $\sin\phi\ge\phi/2$ for small positive $\phi$, the term $\cos U\sin\phi$ controls the sign even at the first band. The other two terms have combined magnitude $O(\phi\sigma^n)$. Thus the next row has sign $(-1)^h$ at the old zero. Taking $\eta<\pi/12$ ensures that the new phase at any old zero in $0<\phi<\eta$ is within $\pi/3$ of the same center $h\pi$. Its sign and the uniqueness and orientation of the new crossing put the new $h$th zero before the old $h$th zero, which is before the new $(h+1)$st zero. The larger interval $0\le\phi\le2\eta$ provides the final comparison for large $n$. This proves strict consecutive alternation, including an initial number of zeros proportional to $n$.

Finally fixed $h$ gives $(n+1)\phi\to h\pi$ by the residue estimate, bounded index, and $B(0)=0$. Substituting into $q(\phi)$ gives the constant. $\square$

\section{The microscopic limit and derivatives}

The preceding two-pole calculation gives a locally uniform limit in a complex neighborhood after rescaling. In fact, uniformly on every compact subset of the entire $x$-plane,
\[
\frac{O_{n,j}(-q_*-x/n^2)}{O_{n,j}(-q_*)}
 \longrightarrow S_\kappa(x)
 =\frac{\sin(\kappa\sqrt{x})}{\kappa\sqrt{x}}
 ={}_0F_1\!\left(;\tfrac32;-\frac{\kappa^2x}{4}\right),
\]
with the value at zero defined by continuity. To see this first write $x=y^2$. For bounded complex $y$, $\phi=\kappa y/n+O(n^{-3})$, $r(\phi)^{n+1}/\rho^{n+1}\to1$ and $a(\phi)\to a(0)$. The normalized error divided by $\phi$ is exponentially small. The sine divided by $\phi(n+1)$ converges uniformly to $\sin(\kappa y)/(\kappa y)$, including zero, by taking the analytic extension of $\sin u/u$. In particular the value of $f_n$ at the fold is asymptotic to $2(n+1)\rho^{-(n+1)}/a(0)$. Dividing by this exact positive value proves the stated convergence; both sides are even in $y$.

For every fixed integer $h\ge0$, all nonzero $h$th row derivatives are positive at every point of $[-q_*,0]$. Indeed their exponential generating function at $-q$ is $h!C(z)^h/D(z,q)^{h+1}$. The critical positivity and its subcritical extension imply strictly positive coefficients whenever $n\ge jh$. For $q>q_*$ this generating function is analytic on the positive real axis, tends to zero at positive infinity, and is nonpolynomial. The Pringsheim argument of the sharp-edge theorem therefore gives infinitely many rows with negative $h$th derivative at $-q$.

Locally uniform convergence gives convergence of every fixed derivative of the rescaled polynomial. Since
\[
S_\kappa^{(h)}(x)=
 \frac{(-\kappa^2/4)^h}{(3/2)_h}\,
 {}_0F_1\!\left(;h+\tfrac32;-\frac{\kappa^2 x}{4}\right),
\]
its positive zeros are $b_{h+1/2,k}^2/\kappa^2$, where $b_{\nu,k}$ is the $k$th positive zero of the Bessel function $J_\nu$. These zeros are simple. For each fixed $k$, apply uniform convergence on small disks about the first $k$ zeros and on their intervening real compact intervals. Rouch\'{e}'s theorem, conjugation, and the absence of earlier real zeros below the fold identify the first $k$ negative real roots of the row derivative. If $q_{n,h,k}$ is the $k$th smallest positive modulus among these real roots, then
\[
\lim_{n\to\infty} n^2(q_{n,h,k}-q_*)=\frac{b_{h+1/2,k}^2}{\kappa^2}.
\]
For $h=0$ this recovers $b_{1/2,k}=k\pi$. See~\cite[\S\S10.2, 10.21]{dlmf}.

\begin{thebibliography}{9}
\bibitem{original} P. Alexandersson, \emph{Real-rooted polynomials: open problems}, The symmetric functions catalog, \href{https://www.symmetricfunctions.com/realRootedOpenProblems.htm\#openProblemsSetPartitions}{section on set partitions and trees}, accessed November 25, 2026.
\bibitem{companion} ProgramComputer, \emph{Three-pole interlacing for ordered partitions and the largest zeros at arbitrary thresholds}, \href{https://github.com/ProgramComputer/open-math-research/blob/main/topics/ordered-large-blocks/threshold-three.pdf}{companion manuscript} (2026).
\bibitem{dlmf} NIST Digital Library of Mathematical Functions, \href{https://dlmf.nist.gov/10}{Chapter 10: Bessel functions}, sections 10.2 and 10.21.
\end{thebibliography}
\end{document}
